%% 04-evaluation.tex --- included by ../sfc.tex
\section{Experimental Evaluation}
\label{sec:experiments}

\subsection{Implementation and Verification}
\label{sec:perf}

\PORTER{} is implemented in Python and is publicly available together
with the verifier and comparison harness. It takes no seed: heuristic
ties resolve through a fixed chain of deterministic keys and any
residual tie through a constant-initialised generator, and every
residual tie-break is counted in the emitted statistics, so a single
routing pass on a fixed circuit and hardware emits a byte-identical
operation stream across runs and behaviour preservation across
refactorings is certified by hashing it. One caveat applies to the
selection protocol of Section~\ref{sec:horizon} and we state it
plainly: the speculative horizons are guarded by a wall-clock budget, so
on a loaded machine a candidate may be abandoned that would otherwise
have won, and the protocol's output is reproducible only up to that
budget. Each individual route remains exactly reproducible, and the
horizon a run selected is recorded alongside its result. Compile time is
kept low by incremental data structures --- chiefly a copy-on-write DAG
frontier simulator\label{sec:sfc-simstate} shared by extended-set
collection and teleport lookahead, whose snapshots cost
$O(|\text{frontier}|)$ rather than $O(|\text{remaining circuit}|)$ ---
together with cached distance tables, $O(1)$ free-slot accounting, and
apply/restore SWAP scoring in place of layout cloning.

Every reported run is certified by an independent replay verifier that
treats the compiler as a black box: from the emitted operation stream,
the original circuit and the initial layout alone it re-derives gate
coverage (each gate exactly once), per-qubit dependency order,
coupling-edge legality of every local operation, port rules for every
teleport, telegate and packet, and consistency of every reported
statistic. The verifier is itself tested by fault injection --- ten
classes of mutation applied to passing streams, every mutant rejected
(\appVerification{}). All \PORTER{} numbers in this paper are
verifier-certified.

\subsection{Setup}
\label{sec:setup}

\textbf{Architectures.} We use the two grid-of-grids topologies
\dSABRE{}~\cite{lidsabre} evaluates on: architecture~A, a $2{\times}2$
array of $3{\times}3$ cores (36 physical qubits, four links), and
architecture~H, a $2{\times}3$ array of $4{\times}4$ cores (96
physical, seven links). Sections~\ref{sec:scaleup}
and~\ref{sec:heavyhex} add its $5{\times}5$-core scalability series and
two networks of IBM 27-qubit heavy-hex tiles. All are read from the
baselines' own device files, so every compiler sees identical physical
qubit ids.

\textbf{Benchmarks.} We use the same MQT-Bench~\cite{quetschlich2023mqtbench}
circuit files \dSABRE{}'s own evaluation uses, confirmed by exact
CX-count match against its published tables: a 25-qubit suite (AE, GHZ,
Graphstate, QFT, QNN, Random) on architecture~A and a 64-qubit suite
(the same six families, plus QPEexact, QAOA and Multiplier) on
architecture~H, with QNN the deeper $8{,}126$-CX variant. To test at
intermediate scale on the same architecture we add an 80-qubit suite of
the same nine families. The 80-qubit AE and QNN instances cannot be
produced by the stock generator and are reconstructed to match its own
construction exactly --- in particular QNN is the deep $12{,}718$-CX
ZZ-feature-map circuit, not the degenerate 79-CX chain the current
generator emits, whose interaction graph is identical to GHZ's.
\appSuite{} details both reconstructions, their byte-level validation,
and the CX counts.

\textbf{Baselines.} We compare against
\TeleSABRE{}~\cite{russo2025telesabre} and \dSABRE{}~\cite{lidsabre},
the latter in its default \texttt{dSABRE\_BurstExt} mode with its
shipped configuration --- which since its most recent revision means
\emph{safe mode}: capacity as a legality condition with a per-core
reserve, BFS slack relay, and a guaranteed per-gate recovery
transaction (Section~\ref{sec:related}). Its router has changed
materially since the version its published tables report, so every
\dSABRE{} number here was regenerated from revision \texttt{cb5515b}
(2026-08-21) rather than quoted, and the harness that produced them is
released with \PORTER{}.

Three protocol decisions keep that comparison honest, and all three
favour the baseline. We run \dSABRE{} under its own budget convention,
letting all three of its recovery budgets derive from the architecture
as its own drivers do rather than imposing values of our own --- a
binding iteration budget would cut its heuristic search short and hand
the remainder to its guaranteed drain at a worse EPR count. We do not charge it for pre-routing layout repair: in safe mode
\texttt{route} brings an externally supplied layout up to its reserve
before the clock starts, and those relay teleports land in its EPR
counter although they adjust the \emph{starting layout} rather than
route anything; \PORTER{} performs the same repair as a free
permutation, so we net it out uniformly. And we
hand it both families of its own layout candidates rather than picking
one (Section~\ref{sec:protocols}), which can only strengthen it
relative to its own published protocol.

\TeleSABRE{}'s stock reference C binary crashes deterministically on
architecture~A, in its Hungarian initial-layout pass rather than in
routing: a hard-coded capacity threshold leaves some qubits unplaced and
their sentinel index is then dereferenced. All \TeleSABRE{} numbers here
use a minimally patched binary that falls back to the least-loaded core;
the patch is in our repository. Every compiler invocation is capped at $900$\,s
wall-clock, uniformly across all three compilers, and a call exceeding
it is recorded like any other failure to complete.

\subsection{Initial-Layout Protocols}
\label{sec:protocols}

We evaluate under three initial-layout protocols, designed to separate
routing quality from initial-layout quality.

\textbf{Own-protocol comparison.} Each router is first evaluated under
its own established best-practice protocol. For the two baselines that
is best-of-three, matching the convention already used to compare
\TeleSABRE{} and \dSABRE{} to each other~\cite{lidsabre}:
\TeleSABRE{} runs its built-in Hungarian initial-layout heuristic at
config seeds 0--2, and \dSABRE{} runs Qiskit \texttt{SabreLayout} on
its corners-removed architecture graph at seeds 0--2, each layout
routed by \SABRE{}-style forward\,$\to$\,backward\,$\to$\,forward
passes; both report the best EPR count. \PORTER{} instead builds one
deterministic temporal locality-greedy layout
(Section~\ref{sec:sfc-layout}) and routes it at each of the three
candidate lookahead horizons of Section~\ref{sec:horizon}, keeping the
best --- three routing passes over one layout, against nine to eighteen
over three to six. It takes no seed and searches no layouts, so every
number is exactly reproducible from the circuit file alone.

\dSABRE{}'s corner removal reserves low-degree physical qubits as
communication slots by excluding them from the SabreLayout coupling map.
It ships two variants of the rule and neither dominates on our
architectures, so rather than choose we give \dSABRE{} both families and
take its best: up to six candidate layouts instead of three. This can
only strengthen the baseline relative to its own published protocol, and
it is why our \dSABRE{} numbers sit at or below the figures it publishes
for the architectures it reports.

\textbf{Shared-layout comparison.} To separate layout quality from
routing quality, both \dSABRE{} and \PORTER{} are then run from an
identical externally supplied mapping --- \TeleSABRE{}'s own initial
layout, and separately \texttt{pytket-dqc}'s hypergraph partition
(\appPytket{}) --- so the only difference left is the routing loop.
\textbf{Cross-layout comparison.} Finally each compiler routes the
other's own layout, which closes the remaining cell of the design.

\subsection{Own-Protocol Results}
\label{sec:main_results}

\TeleSABRE{}'s EPR count is teledata plus telegate operations;
\dSABRE{}'s is teledata only; \PORTER{}'s is teleport plus telegate
plus telegate-packet operations (one EPR pair per packet, regardless
of how many remote CX gates it amortises). All three report in the
same one-EPR-pair unit. SWAP counts intra-core SWAPs only.


\begin{table}[t]
\centering
\caption{Main comparison, own protocol: each compiler from its own
  initial layout. \TeleSABRE{} and \dSABRE{} report the best of three
  seeds; \PORTER{} routes one deterministic layout at each of three
  candidate lookahead horizons (Section~\ref{sec:horizon}), and
  \PORTER{}+R spends five routes on the backward-refinement schedule of
  Section~\ref{sec:sfc-refine} instead --- closer to \dSABRE{}'s nine to
  eighteen, though still short of them. Richer schedules that do reach
  that budget, and pass \dSABRE{}, are in Table~\ref{tab:protocols}.
  $\Delta_\mathrm{TS}$ is the EPR change of whichever \PORTER{} column is
  stronger on that row, against \TeleSABRE{} (\%); negative is better.
  The bolded \PORTER{} column is that stronger one --- refinement leads at
  25 and 64 qubits, horizon selection alone at 80. $^{\ast}$\TeleSABRE{} does not converge in
  three seeds; each gmean row covers the circuits it completes. Local
  SWAP counts, which follow the same ordering throughout, are tabulated
  in \appTables{}.}
\label{tab:main}
\setlength{\tabcolsep}{4pt}
\small
\resizebox{\columnwidth}{!}{%
\begin{tabular}{@{}ll rrrr r@{}}
\toprule
Suite & Circuit & \TeleSABRE{} & \dSABRE{} & \PORTER{} & \PORTER{}+R & $\Delta_\mathrm{TS}$ \\
\midrule
25q & ae          &        119 &    65 &    65 &    65 &   $-45.4$ \\
25q & ghz         &         13 &     4 &     4 &     4 &   $-69.2$ \\
25q & graphstate  &         21 &     6 &     7 &     4 &   $-81.0$ \\
25q & qft         &        138 &    81 &    74 &    74 &   $-46.4$ \\
25q & qnn         &        515 &   143 &   119 &   116 &   $-77.5$ \\
25q & random      &        636 &   404 &   401 &   383 &   $-39.8$ \\
\multicolumn{2}{@{}l}{\quad\textbf{gmean} (6)} & 106.6 & 44.0 & 43.1 & \textbf{38.8} & $\mathbf{-63.6}$ \\
\midrule
64q & ae          &        353 &   165 &   157 &   177 &   $-55.5$ \\
64q & ghz         &         16 &     6 &     6 &     6 &   $-62.5$ \\
64q & graphstate  &         62 &    15 &    20 &    16 &   $-74.2$ \\
64q & qft         &        312 &   179 &   262 &   262 &   $-16.0$ \\
64q & qnn         &        889 &   538 &   638 &   537 &   $-39.6$ \\
64q & random      &  $^{\ast}$ &   732 &   737 &   737 &       --- \\
64q & qpeexact    &        289 &   187 &   222 &   222 &   $-23.2$ \\
64q & qaoa        &        956 &   555 &   775 &   607 &   $-36.5$ \\
64q & multiplier  &       1601 &  1285 &  1076 &   763 &   $-52.3$ \\
\multicolumn{2}{@{}l}{\quad\textbf{gmean} (8)} & 284.5 & 144.6 & 166.2 & \textbf{149.2} & $\mathbf{-47.6}$ \\
\midrule
80q & ae          &  $^{\ast}$ &   332 &   299 &   490 &       --- \\
80q & ghz         &         26 &     7 &     7 &     7 &   $-73.1$ \\
80q & graphstate  &         87 &    20 &    28 &    24 &   $-72.4$ \\
80q & qft         &  $^{\ast}$ &   315 &   446 &   393 &       --- \\
80q & qnn         &  $^{\ast}$ &   791 &  1070 &  1181 &       --- \\
80q & random      &  $^{\ast}$ & 10551 & 11040 & 11040 &       --- \\
80q & qpeexact    &        784 &   347 &   291 &   526 &   $-62.9$ \\
80q & qaoa        &  $^{\ast}$ &   862 &  1122 &  1122 &       --- \\
80q & multiplier  &  $^{\ast}$ &  3793 &  2411 &  2699 &       --- \\
\multicolumn{2}{@{}l}{\quad\textbf{gmean} (3)} & 121.0 & 36.5 & \textbf{38.5} & 44.5 & $\mathbf{-68.2}$ \\
\bottomrule
\end{tabular}%
}
\end{table}



Against \TeleSABRE{} the result is one-sided. From one deterministic
layout \PORTER{} lowers geometric-mean EPR by $59.6\%$, $41.6\%$ and
$68.2\%$ at 25, 64 and 80 qubits, spending three routing passes against
\TeleSABRE{}'s three seeded attempts. Refinement adds to that at 25 and
64 qubits ($63.6\%$ and $47.6\%$) but not at 80, where the single
selected pass is already the stronger of the two ($68.2\%$ against
$63.2\%$) --- the first place in this evaluation where choosing the
lookahead horizon (Section~\ref{sec:horizon}) beats spending four more
routing passes. Every converged circuit improves, from $-16.0\%$ to
$-81.0\%$. \PORTER{} completes all
twenty-four instances, as does \dSABRE{}; \TeleSABRE{} completes
seventeen, failing on 64q~Random and on six of the nine 80-qubit
circuits including the $23{,}381$-CX Random.

Two details are worth recording. On all eight 64-qubit circuits
\TeleSABRE{} completes, our harness reproduces the EPR and SWAP counts
\dSABRE{} publishes for it~\cite{lidsabre} exactly, so the two
evaluations share a baseline \emph{measurement}, not merely a baseline
tool. And the
chain-structured GHZ that defeats locally-greedy schedulers, where
\TeleSABRE{} spends 16 pairs on a 63-CX circuit, costs \PORTER{} and
\dSABRE{} 6 each: \PORTER{}'s first-appearance visit order puts
successive chain neighbours in one core, dissolving the pathology
before routing begins.

Against \dSABRE{} the answer depends on the pass budget and we report
every point. At its own three-route protocol \PORTER{} takes the
geometric mean at 25 qubits ($43.1$ against $44.0$) and trails at 64 and
80 ($166.2$ against $144.6$; $38.5$ against $36.5$, both on the
\TeleSABRE{}-matched basis), winning eight instances, losing twelve and
tying four. Refinement comes within $3.2\%$ at 64; the $\gamma$-search
of Section~\ref{sec:layoutsearch} closes the 64-qubit suite outright
($172.7$ against $173.1$ over all nine circuits) and leaves
$1.08\times$ at 80. Where \PORTER{} wins it wins large ---
64q~Multiplier is $763$ against $1{,}285$ with refinement, the widest
margin here --- and where it loses worst the gap is $1.5\times$, on
80q~QPEexact ($526$ against $347$) and 64q~QFT ($262$ against $179$);
the first of those is a loss for refinement alone, the selected route
taking the instance outright at $291$. This is a comparison against a
target that moved: the \dSABRE{} revision measured here lowers its EPR
on every one of our instances relative to its published tables, by up to
$40\%$ (Section~\ref{sec:related}). Against that earlier version a
single \PORTER{} pass led on all three suites; against this one it does
not, and we report the current numbers.

\subsection{Shared-Layout Results}
\label{sec:shared_results}

The own-protocol comparison mixes two effects: the quality of each
compiler's initial layout and the quality of its routing loop.
Table~\ref{tab:shared} removes the first --- both \dSABRE{} and
\PORTER{} route, single-shot, from the identical layout \TeleSABRE{}
built for itself --- leaving a pure routing comparison.

\begin{table}[t]
\centering
\caption{Shared-layout comparison: \dSABRE{} and \PORTER{} both routed
  from the identical layout \TeleSABRE{} built for itself, so the only
  difference left is the routing loop. Geometric-mean EPR
  over the circuits \TeleSABRE{} completes and therefore provides a
  layout for --- all six at 25q, eight of nine at 64q, three of nine at
  80q. $\Delta_\mathrm{TS}$ columns are each compiler's change against
  \TeleSABRE{}'s own-protocol result. Per-circuit rows are in
  \appTables{}.}
\label{tab:shared}
\setlength{\tabcolsep}{4pt}
\small
\begin{tabular}{@{}l r rr rr@{}}
\toprule
 & & \multicolumn{2}{c}{\dSABRE{}} & \multicolumn{2}{c}{\PORTER{}} \\
\cmidrule(lr){3-4}\cmidrule(lr){5-6}
Suite & \TeleSABRE{} & EPR & $\Delta_\mathrm{TS}$ & $+$R & $\Delta_\mathrm{TS}$ \\
\midrule
25q (6) & 106.6 & 51.0 & $-52.2$ & \textbf{47.7} & $\mathbf{-55.2}$ \\
64q (8) & 284.5 & 189.2 & $-33.5$ & \textbf{157.2} & $\mathbf{-44.7}$ \\
80q (3) & 121.0 & 74.2 & $-38.7$ & \textbf{55.7} & $\mathbf{-53.9}$ \\
\textbf{all} (17) & 173.0 & 100.9 & $-41.7$ & \textbf{86.0} & $\mathbf{-50.3}$ \\
\bottomrule
\end{tabular}
\end{table}


This is the cleanest comparison in the paper, and it is the one
\PORTER{}'s own protocol loses. At three routes each \dSABRE{} leads at
25 and 80 qubits and \PORTER{} at 64 ($-52.2\%$ against $-41.5\%$,
$-38.7\%$ against $-33.1\%$, $-33.5\%$ against $-36.5\%$), pooling to
$100.9$ against $107.7$ over the seventeen instances (the selected
column is in \appTables{}; Table~\ref{tab:shared} reports the refined
one). Spending those three routes on lookahead horizons is therefore
not what closes the gap; spending five on the
forward--backward--forward schedule is, and by a clear margin ---
\PORTER{}+R reaches $86.0$ against \dSABRE{}'s $100.9$ and leads at
every scale. So it is that schedule, not the plan-committed loop, that
separates the two compilers on fixed layouts;
Section~\ref{sec:ablation} confirms this from the other direction:
removing plan commitment moves EPR by only $0.6\%$, and in \PORTER{}'s
favour at that.

\textbf{Pooled fixed-layout grid.} A second source of externally
supplied layouts --- \texttt{pytket-dqc}'s hypergraph placements,
routed by both compilers on the 25q and 64q suites (\appPytket{};
\appOptin{} offers that partition to \PORTER{} as an extra layout
candidate instead) --- pools with Table~\ref{tab:shared} to give
forty-one instances on which \dSABRE{} and \PORTER{} routed the
\emph{identical} layout. Under \PORTER{}'s own protocol \dSABRE{} wins
twenty to sixteen with five ties, pooled geometric mean $134.2$ against
$139.0$; with the refinement schedule that matches its pass count
\PORTER{} wins twenty to fifteen with six ties, $126.6$ against
$134.2$. Routing quality, with the layout controlled and the pass
budget matched, does not separate these two compilers.

\textbf{Each compiler on the other's layout.} Closing the remaining
cell --- routing each system from its own and from its counterpart's
layout (\appTables{}) --- gives the same reading from the other side.
Down a column the compilers are close: on \dSABRE{}'s layout,
\dSABRE{} reaches $163.9$ and \PORTER{}+R $161.1$; on \PORTER{}'s,
$171.2$ and $171.4$. Across a row the layouts are not: \dSABRE{}'s
candidate set is the better start for \emph{both}, by $4.3\%$ for its
own router and $6.0\%$ for \PORTER{}'s. The best cell is \PORTER{}+R on
\dSABRE{}'s layout.

At the sizes this grid covers, then, the routing stages are worth about
the same given the same start, and what separates the two end-to-end is
that \dSABRE{} searches up to six candidate layouts where \PORTER{}
builds one --- a deliberate trade (Section~\ref{sec:time}) that
Section~\ref{sec:layoutsearch} prices. It does not hold at every size:
Section~\ref{sec:scaleup} shows the same control becoming
counter-productive from $150$ physical qubits upward, where a different
parameter binds.

\subsection{Scaling to 360 Logical Qubits}
\label{sec:scaleup}

\begin{table}[t]
\centering
\caption{Scale-up suite: QFT and QPEexact at 100, 200, and 360 logical
  qubits, EPR pairs. Architectures follow \dSABRE{}'s own scalability
  protocol --- $5{\times}5$ cores throughout, core count $6 \to 12 \to
  20$ --- read from the baselines' own device files so all routers see
  identical physical qubit ids; our \dSABRE{} column reproduces the
  figures it publishes for exactly this series~\cite{lidsabre}
  ($201/246/766/896/1890/1938$) to within the one candidate layout our
  protocol adds. The last two columns are the fixed-layout control:
  \PORTER{} routing \dSABRE{}'s own \texttt{SabreLayout} candidates.
  ``$t$'' is one \dSABRE{} attempt (seconds); its protocol runs three
  passes over each of up to six candidates. \texttt{pytket-dqc} is
  omitted: its hypergraph partitioner does not finish in reasonable time
  at these widths.}
\label{tab:scaleup}
\setlength{\tabcolsep}{3.5pt}
\small
\resizebox{\columnwidth}{!}{%
\begin{tabular}{@{}ll rrr rr rr rr@{}}
\toprule
 & & \multicolumn{3}{c}{architecture} & \multicolumn{2}{c}{\dSABRE{}}
 & \multicolumn{2}{c}{\PORTER{}, own layout}
 & \multicolumn{2}{c}{\PORTER{}, \dSABRE{} layout} \\
\cmidrule(lr){3-5}\cmidrule(lr){6-7}\cmidrule(lr){8-9}\cmidrule(lr){10-11}
Circuit & CX & $n$ & $P$ & $K$ & EPR & $t$ & 1 pass & $+$R & 1 pass & $+$R \\
\midrule
QFT      &   3{,}420 & 100 & 150 &  6 &  200 &  1.7 &  198 &  339 &  240 &  273 \\
QPEexact &   3{,}579 & 100 & 150 &  6 &  246 &  2.0 &  200 &  262 &  200 &  283 \\
QFT      &   7{,}220 & 200 & 300 & 12 &  765 &  6.5 &  579 & 1128 & 1031 & 1018 \\
QPEexact &   7{,}579 & 200 & 300 & 12 &  896 &  7.8 &  528 &  979 &  651 &  946 \\
QFT      &  13{,}300 & 360 & 500 & 20 & 1890 & 29.9 & 1156 & 2822 & 1870 & 2546 \\
QPEexact &  13{,}979 & 360 & 500 & 20 & 1938 & 38.9 & 1224 & 2342 & 1444 & 2416 \\
\midrule
\textbf{gmean} (6) & & & & & \textbf{705.7} & & \textbf{507.7} & 930.3 & 665.7 & 877.9 \\
\bottomrule
\end{tabular}%
}
\end{table}

Every result so far is at 108 physical qubits or fewer.
Table~\ref{tab:scaleup} takes the two circuit families that exist at
100, 200 and 360 logical qubits onto \dSABRE{}'s own scalability
series: $5{\times}5$ cores throughout with the core count growing
$6 \to 12 \to 20$, so occupancy is roughly constant while the core
graph grows. It is a monotone ladder in the quantity that matters: inter-core
traffic scales with core count, and \dSABRE{}'s own cost rises
$200 \to 765 \to 1{,}890$ on QFT across it.

\PORTER{} takes all six instances, and by a widening margin: $0.90\times$
\dSABRE{}'s geometric-mean EPR at 100 logical qubits, $0.67\times$ at 200
and $0.62\times$ at 360, for $0.72\times$ overall. It does so at
$0.20\times$ \dSABRE{}'s complete compile cost ($8.4$\,s against
$41.1$\,s in geometric mean), because \dSABRE{} routes up to eighteen
passes here against \PORTER{}'s three.

\textbf{What carries that result is the lookahead horizon, and nothing
else does.} Both mechanisms that close \PORTER{}'s gap below $108$
physical qubits are inert or harmful at these widths. Backward
refinement loses on all six instances, by $83\%$ in geometric mean.
Handing \PORTER{} \dSABRE{}'s own \texttt{SabreLayout} candidates ---
the fixed-layout control that reaches parity at 25 and 64 qubits
(Section~\ref{sec:shared_results}) and \emph{wins} on both heavy-hex
core graphs (Section~\ref{sec:heavyhex}) --- makes it \emph{worse} here,
$0.94\times$ against $0.72\times$, and worst on the largest instance,
where \dSABRE{}'s layout costs $1{,}870$ pairs against $1{,}156$ from
\PORTER{}'s own. The mapping is not the binding constraint at this
scale. The horizon is: every instance in this series selects
$L_\mathrm{tp}{=}128$, the longest candidate, and at the shipped $16$
\PORTER{} instead trails \dSABRE{} by $1.45\times$
(Section~\ref{sec:horizon}).

This separates two deficits the smaller suites conflate. Up to $108$
physical qubits what limits \PORTER{} is its initial layout: control it
and match the pass budget and the difference between the two compilers
disappears. From $150$ upward the layout is no longer the constraint and
the horizon is --- consistent with its being an absolute gate count that
does not grow with $K$, so that at twenty cores a plan is scored over
proportionally far less of the circuit than at four. Selecting it per
circuit is a workaround rather than a fix, and a $K$-aware horizon
remains open.

\subsection{Architecture Independence: Heavy-Hex Cores}
\label{sec:heavyhex}

\begin{table}[t]
\centering
\caption{Architecture-independence check: the nine-circuit 64-qubit
  suite on four IBM 27-qubit heavy-hex cores (108 physical qubits),
  wired as a \emph{ring} (4 links, every core degree 2) and a
  \emph{star} (3 links, core~0 a hub every route must cross).
  Geometric-mean EPR; the right-hand pair is the fixed-layout control,
  \PORTER{} routing \dSABRE{}'s own \texttt{SabreLayout} candidates.
  \PORTER{} completes all eighteen instances, so every row is over the
  full nine; \PORTER{}+R, which does not select a lookahead horizon,
  still fails star Multiplier and is therefore ``---'' on that row.
  Per-circuit rows are in \appTables{}. \dSABRE{} reproduces its published numbers
  for these two architectures on all eighteen instances~\cite{lidsabre}.}
\label{tab:heavyhex}
\setlength{\tabcolsep}{4pt}
\small
\begin{tabular}{@{}l r rr rr@{}}
\toprule
 & & \multicolumn{2}{c}{own layout} & \multicolumn{2}{c}{\dSABRE{} layout} \\
\cmidrule(lr){3-4}\cmidrule(lr){5-6}
Core graph & \dSABRE{} & \PORTER{} & $+$R & \PORTER{} & $+$R \\
\midrule
ring (9)          &  85.2 &  86.1 & 82.3 &  82.9 & \textbf{81.4} \\
star (9)          & 101.1 & 122.3 &  --- &  98.5 & \textbf{93.5} \\
\midrule
star Multiplier   &   722 & \textbf{455} &  --- &   462 &   462 \\
\bottomrule
\end{tabular}
\end{table}


Every result so far uses square grid cores wired in a grid. To test how
much of \PORTER{}'s behaviour depends on that we adopt \dSABRE{}'s own
architecture-independence experiment~\cite{lidsabre}: the same
nine-circuit 64-qubit suite on four IBM 27-qubit heavy-hex tiles. Two
things change at once. The \emph{core} becomes irregular --- degrees
over $\{1,2,3\}$, six degree-1 leaves, diameter 12 against a
$4{\times}4$ grid's 6 --- and the \emph{core graph} is either a ring,
where every core has degree 2, or a star, where core~0 is a hub and the
others are leaves. \PORTER{}'s hardware object is derived from the
architecture \dSABRE{}'s own builder produces, and a topology assertion
checks both compilers see identical qubit ids, intra-core edges and
port endpoints before anything is routed; \dSABRE{} routes on its
native object, reproducing its published numbers on all eighteen
instances. Nothing in either compiler is specialised for a heavy-hex
tile.

\textbf{The core graph matters; the core shape does not.}
Table~\ref{tab:heavyhex} shows an irregular
tile costs \PORTER{} nothing: on the ring \PORTER{}+R \emph{leads}
\dSABRE{} ($82.3$ against $85.2$), winning six of nine and tying GHZ ---
better relative standing than it holds on the grid suite of
Table~\ref{tab:main}. Doubling the intra-core diameter and replacing
degree-4 grid interiors with degree-1 leaves disturbs neither compiler.
The star does: under its own protocol \PORTER{} trails there, $122.3$
against $101.1$, winning two of nine, and \PORTER{}+R --- which cannot
route Multiplier at all --- trails by $13.6\%$ on the eight it
completes.

\textbf{Star Multiplier is a livelock the horizon escapes.} At a fixed
$L_\mathrm{tp}{=}16$ this instance does not terminate, and it fails as a
local-SWAP livelock rather than the capacity deadlock of
Section~\ref{sec:sfc-progress}: retired gates freeze at $3{,}226$ of
$13{,}040$ and the EPR count at $538$ while intra-core SWAPs accumulate
at one per iteration, so the iteration guard rather than the progress
argument stops it. It is a property of the horizon, not of the circuit:
the same layout and the same router complete it at
$L_\mathrm{tp}{=}128$, at $455$ pairs against \dSABRE{}'s $722$. Since
horizon selection routes that candidate as a matter of course,
\PORTER{} has no unrouted instance anywhere in this evaluation ---
whereas \PORTER{}+R, which holds the horizon fixed, still fails it.

\textbf{The star gap is the layout, not the routing.} Given
\dSABRE{}'s own layout \PORTER{} completes all nine star circuits ---
including Multiplier at $462$ pairs against \dSABRE{}'s $722$, which its
own layout now also routes, slightly better, at $455$ --- and
with the pass budget matched the control reverses the result on
\emph{both} core graphs ($-4.4\%$ on the ring, $-7.5\%$ on the star),
the larger margin falling on the topology \PORTER{} loses outright
under its own protocol. The cause is visible in the layout:
\texttt{build\_initial\_layout} reserves $f_\mathrm{min}$ in
\emph{every} core, treating cores as interchangeable. They are on a
grid and on a ring, where every core has the same degree; they are not
on a star, where Multiplier opens with free counts $(2,2,13,27)$ ---
the hub every route must cross pinned at the floor while a whole core
sits empty. Raising the reservation uniformly does not repair it (we
tried $f_\mathrm{min}{\in}\{4,6\}$); what the star asks for is a
reservation scaled by a core's position in the core graph, which we have
not implemented.

Taken with Section~\ref{sec:scaleup} this bounds the layout reading
rather than extending it. Up to $108$ physical qubits --- across two
core shapes and three core graphs --- the initial layout is what
separates the two compilers, and controlling it erases the difference.
Above $150$ the layout stops being the constraint, and
Section~\ref{sec:horizon}'s parameter takes over.

\subsection{Selecting the Lookahead Horizon}
\label{sec:horizon}

The teleport-lookahead horizon $L_\mathrm{tp}$ is the one parameter of
Section~\ref{sec:sfc-scoring} that no single value serves. In an
offline sweep over $\{4,8,\dots,1024\}$ on every instance the best value
spans the whole range: $80$-qubit QNN wants $4$, the $25$-qubit suite
wants $16$, and every $100$-, $200$- and $360$-qubit instance wants
$128$, beyond which EPR saturates. The gap between the shipped $16$ and
the per-instance best is $17.2\%$ geometric-mean EPR.

The direction is set by how dense the circuit's interaction graph is.
Writing $\rho$ for the mean fraction of the other logical qubits each
qubit shares a CX with, a long horizon reports the moving qubit as needed
in \emph{every} core at once when $\rho\to1$, so the gain estimate
degenerates into noise and short horizons win; when $\rho$ is small the
signal is coherent and long horizons win. Across the benchmark, $\rho$
ranks the best horizon at Spearman $-0.765$. It is not, however, a
usable predictor on its own: in the band $0.42\le\rho\le0.53$, where most
$64$-qubit circuits sit, QAOA wants $8$ while AE wants $128$, and a
threshold rule on $\rho$ recovers only $8.0\%$ of the available
$17.2\%$.

\PORTER{} therefore uses $\rho$ to \emph{prune} rather than to choose. A
circuit with $\rho\ge0.80$ is routed at $L_\mathrm{tp}\in\{4,8,16\}$ and
every other circuit at $\{8,16,128\}$; the lowest-EPR route is kept. No
instance above the threshold ever selected the long horizon, so pruning
costs no EPR, and it removes the case where the long horizon fails to
converge and consumes the whole attempt --- the dominant compile-time
term. Two rules keep the cost bounded and the result safe. The
circuit's own default horizon is routed first and unguarded, and wins
ties, so selection can never return a worse route than routing once;
the remaining candidates then share a wall-clock budget proportional to
what that first route cost, and one that overruns it is abandoned. The
budget is what makes a deep horizon affordable to try, and it is also
the one place \PORTER{}'s output depends on the machine it runs on
(Section~\ref{sec:perf}).

Across all $48$ instances this yields $23$ improvements, no regressions,
and $17.4\%$ lower total EPR, at $5.1\times$ the cost of a single pass
--- still below \dSABRE{}'s nine-route protocol
(Section~\ref{sec:time}). It is the mechanism behind the scale-up
result of Section~\ref{sec:scaleup} and behind \PORTER{} completing
every heavy-hex instance in Section~\ref{sec:heavyhex}.

\subsection{Mechanism Ablation}
\label{sec:ablation}

\begin{table}[t]
\centering
\caption{Mechanism ablation: \PORTER{} EPR with one design decision
  disabled at a time, all from the same temporal locality-greedy
  layout. Columns: ``No pk'' disables packets
  (\texttt{max\_packet\_size}=1), ``No rel.'' disables
  both relief tiers, ``Hop-wise'' replaces plan commitment with
  per-iteration re-scoring of the whole front, ``\TeleData{}'' removes
  the \TeleGate{} family; ``pk'' and ``tg'' are the packet and
  single-telegate counts of the full configuration. The last gmean row
  pairs Full against Hop-wise over all twelve rows, which both
  complete. $^{\dagger}$Fails to complete (deadlock
  or iteration-guard trip). $^{\ddagger}$Completes only with the additional
  relief tier described in the text, which the \TeleGate{} otherwise
  makes unnecessary. The first gmean row is over the fourteen instances the
  first three configurations share; the second is the \TeleData{}-only column
  paired against Full over the sixteen instances it completes.}
\label{tab:ablation}
\setlength{\tabcolsep}{3pt}
\small
\resizebox{\columnwidth}{!}{%
\begin{tabular}{@{}ll rrrrr rr@{}}
\toprule
Suite & Circuit & Full & No pk & No rel. & Hop-wise & \TeleData{} & pk & tg \\
\midrule
25q & ae         &    65 &    71 &    80 &    65 & 68$^{\ddagger}$ &   8 &    0 \\
25q & qft        &    74 &    74 &   140 &    74 &   74 &   0 &    0 \\
25q & qnn        &   119 &   122 &   212 &   128 &  119 &   0 &    0 \\
25q & random     &   401 &   412 &   395 &   399 & 413$^{\ddagger}$ &  43 &   93 \\
64q & ae         &   177 &   211 &   327 &   172 & 198$^{\ddagger}$ &  32 &    0 \\
64q & qft        &   262 &   372 &   382 &   336 & \textbf{216}$^{\ddagger}$ &  45 &    1 \\
64q & qnn        &   864 &   931 &   600 &   785 & \textbf{694}$^{\ddagger}$ & 240 &    5 \\
64q & random     &   737 &   843 &   789 &   805 & $^{\dagger}$ &  77 &  228 \\
80q & ae         &   490 &   688 & $^{\dagger}$ &   383 & \textbf{354}$^{\ddagger}$ & 264 &   15 \\
80q & qft        &   446 &   435 & $^{\dagger}$ &   415 & \textbf{307}$^{\ddagger}$ & 209 &    2 \\
80q & qnn        &  1517 &  1571 & $^{\dagger}$ &  1474 & \textbf{926}$^{\ddagger}$ & 448 &   30 \\
80q & random     & 11040 & 11628 & $^{\dagger}$ & 11280 & $^{\dagger}$ & 572 & 3538 \\
\midrule
\multicolumn{2}{@{}l}{\textbf{gmean} (14)} & 57.6 & 61.4 & 67.0 & & & & \\
\multicolumn{2}{@{}l}{\textbf{gmean} (16)} & 78.3 & 83.9 & & & \textbf{71.7} & & \\
\multicolumn{2}{@{}l}{\textbf{gmean} (12)} & 413.3 & & & 411.7 & & & \\
\bottomrule
\end{tabular}%
}
\end{table}

Table~\ref{tab:ablation} attributes \PORTER{}'s EPR consumption to its
mechanisms by disabling each in isolation, holding the initial layout
and every other parameter fixed; the per-action counts behind these
rows are tabulated in \appMix{} (the six omitted rows --- GHZ and
Graphstate at every scale --- form no packets, never trigger relief and
issue no telegates, so all configurations coincide there). Two columns
go further than a per-primitive ablation: ``Hop-wise'' replaces
\PORTER{}'s unit of commitment with \dSABRE{}'s while holding scoring,
plan set and relief tiers fixed, and ``\TeleData{}'' removes the
\TeleGate{} family outright, leaving state teleportation as the sole
inter-core action.

\textbf{Packets.} Disabling \TeleGatePacket{} makes every
shared-control remote CX run pay per gate, and the cost tracks the
packet count: $+42\%$ EPR on 64q~QFT (262 $\to$ 372, 45 packets),
$+40\%$ on 80q~AE (490 $\to$ 688, 264 packets), $+19\%$ on 64q~AE,
against a $+7\%$ overall gmean; the lone counter-move (80q~QFT,
$446 \to 435$) shows packet formation is a heuristic commitment that
can displace a slightly better teleport plan. Read alone, this column
is a clean win for the primitive. The next one overturns that reading.

\textbf{\TeleData{} only.} Removing the \TeleGate{} family exposes two
things the per-primitive ablations cannot. First, the \TeleGate{} is
what makes \PORTER{} \emph{terminate}, not merely what makes it cheap:
it is the only remote action needing no free destination slot, so it is
always a legal plan for an adjacent remote CX. Without it a front whose
remote CXs are all core-adjacent and whose destinations sit at or below
$f_\mathrm{min}$ has an empty plan set, while the tail-relief BFS is
only ever offered the \emph{non-adjacent} blocked candidates --- ten of
the twelve rows deadlock on this. Feeding adjacent-blocked candidates
into the same BFS (marked $^{\ddagger}$, and measurably inert whenever
the \TeleGate{} is available) recovers all but the two Random rows,
which livelock --- and those two are the most telegate-heavy in the
suite ($228$ and $3{,}538$ single telegates), so the primitive is
load-bearing precisely where it is used most.

Second, where \TeleData{}-only completes it is \emph{cheaper than the
full action set}, increasingly so with scale: $+2\%$ at 25q, mixed at
64q, and $-33\%$ gmean at 80q; overall $8.5\%$ below Full on the
sixteen instances it completes. The circuits
it wins are exactly the packet-heaviest, making the ordering there
$\text{\TeleData{} only} < \text{Full} < \text{No packets}$: packets
clearly beat the single telegates they degrade into, but on structured
circuits at scale a pure teleportation policy beats both. We read this
as myopia in the plan comparator rather than a verdict on packets --- a
packet's benefit is realised immediately and in full, a teleport's is
latent and discounted across a bounded horizon --- but we report it as
measured. It also answers a question \dSABRE{} states as open: its
teledata-only vocabulary was chosen by measurement, an ungrouped
telegate serving one gate where teledata amortises many, with grouping
left as the open case. Grouping does close that per-EPR gap against
single telegates --- and still does not justify spending pairs that
teleporting would have spent better.

\textbf{Plan commitment.} Replacing \PORTER{}'s unit of commitment with
\dSABRE{}'s --- returning to the top of the loop after one round of
preparation SWAPs, so the front is re-classified and every plan
re-scored before the next move --- leaves EPR consumption
statistically unchanged: $413.3$ against $411.7$ over the twelve rows
here and $148.8$ against $148.0$ over the twenty-one instances we ran
it on, with eight circuits better without commitment, five worse,
eight tied, and all twenty-one completing either way. The two also make
the \emph{same number of routing moves} ($128{,}222$ against
$128{,}565$ iterations), which is the key number: commitment does not
change what the compiler decides, only what each decision costs. What
it buys is compile time --- $1.31\times$ in geometric mean and
$3.07\times$ on 80q~Random --- because a committed plan is not
re-derived and re-scored every iteration. It is therefore an efficiency
mechanism rather than a quality one: what makes a single \PORTER{} pass
cheap enough to be a competitive operating point, not what makes that
pass good.

\textbf{Capacity relief.} Disabling both relief tiers is not merely
more expensive --- it is \emph{unsafe}: 80q~AE, QFT, QNN and Random all
fail to complete, because with every candidate destination below
$f_\mathrm{min}$ and no mechanism to drain qubits toward reservoirs the
deterministic-greedy loop has no legal remote plan left. Where routing
does complete, dense circuits pay heavily ($+89\%$ on 25q~QFT, $+78\%$
on 25q~QNN, $+85\%$ on 64q~AE). The two instances where relief
\emph{hurts} --- 64q~QNN ($864 \to 600$ without it) and, marginally,
25q~Random ($401 \to 395$) --- show the opportunistic tier can overpay
where congestion would have resolved itself. Of the four mechanisms this table
ablates, capacity relief is the one that earns its place outright: it
converts the raise-on-deadlock failure model of
Section~\ref{sec:sfc-progress} from a liability into a non-event. (The
initial layout, which this table holds fixed rather than ablating,
earns its own place separately --- Sections~\ref{sec:shared_results}
and~\ref{sec:heavyhex}.)

\subsection{Operating Points and Compile Time}
\label{sec:layoutsearch}
\label{sec:time}

\begin{table}[t]
\centering
\caption{Protocol trade-off over all 24 instances. ``routes'' is the number of
  routing passes; times are complete compile costs (geometric mean per
  instance). \PORTER{} routes each of three candidate lookahead
  horizons and keeps the best (Section~\ref{sec:horizon}); \PORTER{}+R adds backward
  refinement (Section~\ref{sec:sfc-refine}) at a fixed horizon;
  \PORTER{}+S adds the best-of-$\gamma$ layout search of this section
  on top of it; \PORTER{}+SR takes the better of horizon selection and
  that search, netting the one route they share. All are seed-free.}
\label{tab:protocols}
\small
\resizebox{\columnwidth}{!}{%
\begin{tabular}{@{}l r rr rr@{}}
\toprule
Protocol & routes & EPR & vs.\ \dSABRE{} & time (s) & vs.\ \dSABRE{} \\
\midrule
\dSABRE{}                 & 9--18 & 163.9 & $1.00$ & 2.26 & $1.00$ \\
\PORTER{}                 &     3 & 174.9 & $1.07$ & 1.32 & $\mathbf{0.58}$ \\
\PORTER{}+R               &     5 & 171.4 & $1.05$ & 1.67 & $0.74$ \\
\PORTER{}+S               & 9--13 & 163.2 & $1.00$ & 4.59 & $2.03$ \\
\PORTER{}+SR              & 11--15 & \textbf{158.2} & $\mathbf{0.97}$ & 5.65 & $2.50$ \\
\bottomrule
\end{tabular}%
}
\end{table}

\PORTER{}'s protocol routes one deterministic layout at each of three
candidate lookahead horizons and keeps the best
(Section~\ref{sec:horizon}), and two further schedules spend more
routes: \PORTER{}+R adds the backward refinement of
Section~\ref{sec:sfc-refine}, and \PORTER{}+S additionally sweeps the
one free scalar in the layout heuristic --- the per-layer discount
$\gamma$, pinned at $0.99$ elsewhere --- over $\{0.90,0.95,0.99,1.0\}$
and keeps the best route by EPR. Both remain deterministic and
seed-free, and the sweep is a search inside \PORTER{}'s own design
rather than an external partitioner or a seed sweep.
A fourth point falls out of the first three. Selection and the
$\gamma$-search look for different things --- the horizon at a fixed
layout, against the layout at a fixed horizon --- and neither dominates:
the search gives the better route on fourteen of the twenty-four
instances, selection on four, and they tie on six. \PORTER{}+SR runs both and keeps the lower EPR,
netting the one route they share. Table~\ref{tab:protocols} places the
four points against \dSABRE{}, pooled over all twenty-four.

The trade is time against EPR pairs, and it is monotone in both.
\PORTER{} costs $1.07\times$ \dSABRE{}'s EPR at $0.58\times$ its compile
time; refinement reaches $1.05\times$ at $0.74\times$; the layout search
reaches parity ($163.2$ against $163.9$) at $2.03\times$; and
\PORTER{}+SR passes it, at $0.97\times$ its EPR for $2.50\times$ its
compile time. Nothing here is free, because \PORTER{} is dearer per
routing pass and the richer schedules run more of them --- $12.4$ on
average for the $\gamma$-search alone.

That last point answers the question Table~\ref{tab:main} raises.
\PORTER{}+SR leads \dSABRE{} on twelve of the twenty-four instances
against nine behind, and by construction never loses to either schedule
it is built from. The residual own-protocol gap is therefore a
\emph{search-budget} gap, not a routing deficit: given a budget in
\dSABRE{}'s range \PORTER{} passes it rather than merely drawing level.
Where on the curve to sit depends on whether EPR pairs or compile
seconds are scarce.

Selecting the horizon has largely overtaken refinement as an
\emph{operating point}: over the twenty-four instances \PORTER{}+R is
ahead by only $2.0\%$ in geometric mean for four extra routes, and it is
$12.0\%$ \emph{behind} the selected pass at 80 qubits, $83\%$ behind on
the scale-up series, and the configuration that fails star~Multiplier.
We report it throughout only because the fixed-layout comparison of
Section~\ref{sec:shared_results} needs a \PORTER{} schedule that spends
its extra effort on routing passes rather than on layouts, which
\PORTER{}+S by construction cannot be.

All times are \emph{complete} compile costs, following \dSABRE{}'s own
convention: layout generation plus every routing pass over every
candidate, not the winning attempt alone. For \dSABRE{} that is nine to
eighteen passes, against three for \PORTER{} and five for \PORTER{}+R;
\TeleSABRE{}'s figure covers all three seeds. \PORTER{}'s advantage is
therefore in the count of passes, not in their price: per routing pass
it is the dearer of the two on all nine 64-qubit circuits, at $0.27\times$
\dSABRE{}'s per-pass speed in geometric mean, partly because enumerating
whole plans, collecting packets by forward simulation and running a
two-tier capacity search all cost more per iteration than scoring a flat
list of hops, and partly because a deep horizon makes each iteration
dearer still. On that nine-circuit subset the complete costs are
$1.98$\,s against $1.86$\,s and \dSABRE{} is the cheaper; pooled over
all twenty-four \PORTER{} is ahead at $0.58\times$.

Over the five circuits \TeleSABRE{} completes (per-circuit times in
\appTables{}), geometric-mean complete cost is $5.95$\,s for
\TeleSABRE{}, $2.67$\,s for \dSABRE{}, $1.55$\,s for \PORTER{} and
$1.97$\,s with refinement, and neither \PORTER{} point holds that margin
on the largest circuit here, where \PORTER{}+R costs $191$\,s against
$144$\,s at $23{,}381$ CX. The cheap operating point is still
\PORTER{}, and at $1.07\times$ \dSABRE{}'s EPR it no longer gives up
much to be so.
